\backhead{FURTHER EXERCISES}{khred}
\begin{exlist}
\item A main of 300 m and 150 mm diameter carries 20 L\,s$^{-1}$ (Hazen--Williams $C=130$). (a) Find the friction loss. (b) The outlet is 4 m higher than the pump and needs 25 m of pressure: find the pump head (add 10 per cent for local losses).\exans{(a) 2.86 m; (b) 32.1 m}
\item The pump of the previous exercise has an efficiency of 0.70. Find the power. How much energy does it use to apply 50 mm to 5 ha?\exans{9.0 kW; 313 kWh}
\item A Kostiakov--Lewis intake $Z=k\tau^a+f_0\tau$ has $k=0.008$ m\,min$^{-a}$, $a=0.4$ and $f_0=0.0001$ m\,min$^{-1}$. Find the depth infiltrated in 60 and 240 min.\exans{47.1 mm; 95.6 mm}
\item Manning: a furrow carries 1.5 L\,s$^{-1}$ on a slope of 0.003 with $n=0.04$. For a trapezoid with bottom 0.1 m and side slope 1, the flow depth is 0.08 m. Find the area, the wetted perimeter, the velocity by Manning and the discharge it gives.\exans{0.0144 m$^2$; 0.326 m; 0.171 m\,s$^{-1}$; 2.5 L\,s$^{-1}$}
\item A field 300 m by 200 m receives a gross depth of 70 mm in a set of 12 h. Find the flow needed. How many sets of 1 ha can a supply of 40 L\,s$^{-1}$ irrigate in 7 days of 24 h, at 70 mm each?\exans{97.2 L\,s$^{-1}$; 34 sets}
\end{exlist}
\backhead{ESSAY QUESTIONS}{khred}
{\small\itshape Answer each question in about one page, with equations, a sketch and at least one number from this chapter.}\par\vspace{3pt}
\begin{exlist}
\item Explain the energy equation for flow in pipes, and how the friction and local losses set the size of a pump.
\item Compare the Kostiakov, Kostiakov--Lewis and Green--Ampt descriptions of infiltration. What does each represent physically, and when is each appropriate?
\item Discuss how weather data are put on the twenty-year grid of SWAFL, and the errors that can arise when a short or incomplete record is used.
\item Explain what satellite vegetation indices can and cannot tell about the water status of a field.
\item Discuss the choice of an irrigation system for a field: which properties of the soil, the water source, the crop and the farm decide between surface, drip and sprinkler?
\item Why are uniformity and efficiency different measures? Show with an example how a uniform system can still be inefficient.
\end{exlist}
\dykband{EXERCISES IN SWAFL}
{\small Open SWAFL (desktop or online) with the benchmark project. The chapter of the \href{https://github.com/halilunaziru73-creator/Coupled-Water-Fluxes-Pore-Scale-Dynamics-Precision-Irrigation-Mulching-Feedback-Loops/raw/main/Software/v1.1.0/SWAFL_Desktop/docs/manual_v1.1.0/SWAFL_Manual_v1.1.0.pdf}{SWAFL manual} that explains each task is given in brackets; record your values in a table and compare them with what is expected (\emph{Check}).}\par\vspace{3pt}
\begin{exlist}
\item {\color{khblue}[Manual, Chapter 20]} Find your farm on the satellite map, demarcate it and read the vegetation and water indices. Comment on the variation within the field.\par\noindent{\small\itshape Check: the indices vary with soil and crop cover.}
\item {\color{khblue}[Manual, Chapter 20]} Add a water source and let the program propose an irrigation system for the field. Record the proposal and the reasons it gives.\par\noindent{\small\itshape Check: the proposal depends on the slope, the soil and the source.}
\item {\color{khblue}[Manual, Chapter 19]} Import a weather file of your own and switch back to the benchmark record. Compare the yearly rain and the reference evapotranspiration.\par\noindent{\small\itshape Check: the comparison shows how your climate differs.}
\item {\color{khblue}[Manual, Chapter 18]} Enter the field size as a length and a width and as an area, and check that the design layers give the same volumes.\par\noindent{\small\itshape Check: the volumes agree.}
\end{exlist}
